\begin{textsample}{POS Dim 4 – human – Score 5.00 – mg/mg\_ef\_linguagens\_lp\_4\_p3.txt}  \label{ex:f4_pos_010}
As práticas pedagógicas de compreensão e produção devem dar preferência a \textbf{textos} reais dos gêneros em circulação na sociedade, começando pelos mais familiares aos estudantes e se encaminhando para os mais distantes de sua experiência imediata.
A função social e estrutura de determinados gêneros é mais facilmente apreensível que a de outros, consideradas as possibilidades da faixa etária e da etapa de desenvolvimento do estudante.
Assim, gêneros do domínio privado, como bilhetes, cartas, convites, são de mais fácil \textbf{leitura} e produção que \textbf{textos} do domínio público, especialmente aqueles produzidos em situações de comunicação formal, como atas ou palestras. É preciso, porém, ter em mente que, ao final da Educação Básica, o estudante deve estar em condições de usar a linguagem \textbf{oral} e escrita em situações públicas de interlocução ( assembleias, palestras, seminários de caráter político, técnico, \textbf{leitura} e produção de \textbf{textos} multissemióticos, de \textbf{textos} científicos, etc. ) e demonstrar disposição e sensibilidade para apreciar os usos artísticos da linguagem.

% matched lemmas: leitura, oral, texto
\end{textsample}
